\section{Full Clifford structures and their tensor maps}\label{cliff:section}

We fix the conventions needed to identify the final correspondence exactly.
The classical construction is reviewed in
\cite[Chapter~4, \S2]{Huybrechts} and \cite[\S\S5--6]{vanGeemen}; we
include the details concerning signs, polarizations, and the two tensor slots.
Our Hodge circle acts on a vector of type $(p,q)$ by
$e^{i(p-q)\theta}$. Tate twists leave this circle unchanged.

Let $H$ be a rational Hodge structure of K3 type: it has weight two,
$h^{2,0}=h^{0,2}=1$, and no Hodge types other than $(2,0)$, $(1,1)$,
and $(0,2)$. Let $Q$ be a
nondegenerate rational quadratic Hodge form. We allow either overall sign:
$Q$ has sign $\epsilon\in\{1,-1\}$ on the real period plane and the
opposite sign on its orthogonal. Write
\[
 C(H,Q)=C^+(H,Q)\oplus C^-(H,Q),\qquad v^2=Q(v,v).
\]
If $H^{2,0}=\mathbb C(e_1+i e_2)$, where $e_1,e_2$ are real orthogonal
vectors with $Q(e_j,e_j)=\epsilon$, set
\begin{equation}\label{cliff:circle}
 J=\epsilon e_1e_2,\qquad
 g_\theta=\cos\theta+J\sin\theta.
\end{equation}
The element $J$ depends only on the oriented period plane. Left
multiplication by $g_\theta$ gives both parity spaces weight-one Hodge
structures, denoted $P_H^+$ and $P_H^-$, with the $i$-eigenspace of $L_J$
as their $(1,0)$-part. We retain these \emph{full} spaces throughout.

\begin{proposition}\label{cliff:construction}
The Hodge structures $P_H^\pm$ are polarizable, and
\begin{equation}\label{cliff:multiplication}
 m^Q:H(1)\longrightarrow\operatorname{Hom}(P_H^+,P_H^-),
 \qquad x\longmapsto L_x
\end{equation}
is a rational Hodge embedding. The parity spaces are rationally
Hodge-isomorphic by right multiplication by any fixed invertible odd
Clifford element. On a sufficiently small period chart, their Hodge
filtrations vary holomorphically and admit fixed rational polarizations.
\end{proposition}

\begin{proof}
The Clifford relations give
\[
 J^2=-1,\quad [J,e_1]=-2e_2,\quad [J,e_2]=2e_1,
 \quad [J,e_1+i e_2]=2i(e_1+i e_2).
\]
Moreover, $J$ commutes with the orthogonal complement of the period
plane. Thus $g_\theta\in\operatorname{Spin}(H,Q)(\mathbb R)$, and its
conjugation action on $H$ is the prescribed weight-two Hodge circle.
Since
\[
 L_{g_\theta}L_xL_{g_\theta}^{-1}
   =L_{g_\theta xg_\theta^{-1}},
\]
the map in \eqref{cliff:multiplication} is a Hodge map. It is injective
because $L_x(1)=x$.
Choose a rational nonisotropic vector $z\in H$. Right multiplication
$R_z$ exchanges the parities, has inverse $R_{z/Q(z,z)}$, and commutes
with the left circle. The same observation applies to any invertible
odd right multiplier.

For holomorphic variation, put $\omega=e_1+i e_2$.
The relation $J\omega=i\omega$ shows that the image of $L_\omega$
from the opposite parity is contained in the $i$-eigenspace of $L_J$.
Conversely $\omega J=-i\omega$, so this eigenspace is contained in
$\ker L_\omega$. Finally
\[
 L_\omega L_{\bar\omega}+L_{\bar\omega}L_\omega
       =4\epsilon\operatorname{id}
\]
shows that $\ker L_\omega=\operatorname{im}L_\omega$, with the
appropriate parities understood. Hence
$F^1P_{H,\mathbb C}^\pm=\operatorname{im}L_\omega$; this is a
holomorphic subbundle as the isotropic line varies.

Here is a fixed rational polarization. Write $u\mapsto u^*$ for
reversal in the Clifford algebra and $\operatorname{scal}$ for its
scalar part. Choose rational orthogonal vectors $f_1,f_2$ of sign
$\epsilon$ whose oriented plane is sufficiently close to the given
period plane. Rational approximation followed by rational
orthogonalization supplies them. Put $a=f_1f_2$ and, on $C^+(H,Q)$,
\begin{equation}\label{cliff:polarizing-form}
 s(u,v)=\operatorname{scal}(uav^*).
\end{equation}
The scalar part is cyclic and invariant under reversal, as is checked
on an orthogonal monomial basis. Since $a^*=-a$, the form $s$ is
alternating. It is nondegenerate because $a$ is invertible and
$(u,v)\mapsto\operatorname{scal}(uv^*)$ is nondegenerate. Reversal
gives $g_\theta^*g_\theta=1$, so $s$ is invariant under the left circle
at every period.

At the reference plane spanned by $f_1,f_2$, write $a=cJ$ with
$c\in\mathbb R^\times$. Then
\[
 s(u,Jv)=c\operatorname{scal}\bigl(\operatorname{Ad}(J)u\,v^*\bigr).
\]
To see definiteness, extend a normalized basis of this plane to an
orthogonal basis with squared lengths
$\epsilon,\epsilon,-\epsilon,\ldots,-\epsilon$.
The last pairing is diagonal on even monomials $e_I$.
If $r=|I|$ and $k$ entries of $I$ belong to the period plane, its
diagonal sign is
\[
 (-1)^k\epsilon^r(-1)^{r-k}=1,
\]
because $r$ is even. Up to the overall sign required by the
weight-one polarization convention, $s$ is therefore a polarization.
Circle invariance makes $s(u,Jv)$ symmetric and nondegenerate at every
period. Its signature is constant on a connected period chart
containing both the reference plane and the original period, so the
chosen sign of $s$ polarizes throughout that chart.
Transport by $R_z$ supplies a fixed polarization
on $C^-(H,Q)$ as well.
\end{proof}

Choose abelian varieties realizing the polarizable weight-one spaces
in Proposition~\ref{cliff:construction}; rational realizations suffice.
For a weight-one Hodge structure $P$, fix the rational Hodge
isomorphism
\[
 j_P:P^\vee(-1)\xrightarrow{\sim}P
\]
induced by its chosen polarization. A map
$f:H(1)\to\operatorname{Hom}(P,M)$ then determines the weight-two
tensor map
\begin{equation}\label{cliff:tensor}
 H\longrightarrow\operatorname{Hom}(P,M)(-1)
   =M\otimes P^\vee(-1)
   \xrightarrow{1\otimes j_P}M\otimes P.
\end{equation}
When $H\subset H^2(S,\mathbb Q)$ for a surface $S$, choose abelian
realizations $M=H^1(A_M,\mathbb Q)$ and $P=H^1(A_P,\mathbb Q)$.
We call $f$ \emph{algebraically realized} if some cycle in
$\operatorname{CH}^2(S\times A_M\times A_P)_{\mathbb Q}$ has, on
$H$, exactly the action \eqref{cliff:tensor} included in the
$H^1\otimes H^1$ K\"unneth component. This specifies the restricted
action on $H$; the algebraic target projectors may be applied to
remove other target components. The terminology always refers to
actual algebraic cycles.

\begin{lemma}[Operations on the two slots]\label{cliff:slots}
Suppose $f:H(1)\to\operatorname{Hom}(P,M)$ is algebraically realized.
For rational weight-one Hodge maps
$a:M\to M'$ and $b:P'\to P$, the map
$x\mapsto a f(x)b$ is algebraically realized. The same holds after
changing the chosen duality isomorphisms, interchanging the tensor
factors, or making rational weight-one Hodge identifications with
isogenous abelian models.
\end{lemma}

\begin{proof}
Under \eqref{cliff:tensor}, the two operations are precisely
\begin{equation}\label{cliff:slot-formula}
 a\otimes\bigl(j_{P'}b^\vee(-1)j_P^{-1}\bigr).
\end{equation}
Both factors are rational weight-one Hodge maps, so
Lemma~\ref{cyc:h1} realizes them by abelian correspondences.
Their product acts independently on the two factors and preserves
the codimension of the correspondence from the surface.
Changing $j_P$ likewise composes with a weight-one Hodge
isomorphism. The factor-switching isomorphism is algebraic; its
K\"unneth sign can be included in the rational coefficient.
The projectors of Lemma~\ref{cyc:projectors} remove all other
K\"unneth components when necessary.
\end{proof}

In particular, the standard full even-Clifford model is $W=P_H^+$.
For a stipulated invertible odd right multiplier $z$, its
endomorphism-valued embedding is exactly
\begin{equation}\label{cliff:standard}
 H(1)\longrightarrow\operatorname{End}(W),\qquad
 x\longmapsto\bigl[w\longmapsto xwz\bigr]
       =R_zL_x.
\end{equation}
Applying the stipulated duality $W^\vee\simeq W(1)$ gives its
specified map to $W\otimes W$. Thus algebraicity of $m^Q$ gives
algebraicity of this precise map, with the fixed multiplier,
polarization, rational normalization, and identification with
$H^1$ of the chosen Kuga--Satake abelian variety. No simple factor
replaces the full even-Clifford space.

These conclusions are also independent of the usual equivalent
Clifford conventions. Reversal takes left multiplication by a spin
element to right multiplication by its inverse, since
$(gu)^*=u^*g^{-1}$. Passing between homology and cohomology uses
duality, transpose, and a weight-one polarization. These are exactly
the operations in Lemma~\ref{cliff:slots}; the source Hodge circle
is always the one fixed above. Reversing the period Hodge action
independently would be a different construction.
